\documentclass[../DoAn.tex]{subfiles}
\begin{document}

Tổng quan chương.

Sau khi đã trình bày phương pháp đề xuất, phân tích cơ sở lý thuyết và kiểm chứng bằng thực nghiệm, đồ án cần khép lại bằng một kết luận trung thực về những gì đã đạt được và những gì vẫn còn bỏ ngỏ. Vì vậy, chương này tổng hợp lại đóng góp chính của đồ án theo đúng trục nghiên cứu, đồng thời nêu rõ giới hạn hiện tại và các hướng phát triển tiếp theo.

Phần đầu của chương nhắc lại kết quả cốt lõi của đồ án ở cả ba mặt: đề xuất lớp điều phối phi tập trung cho MLS, triển khai ứng dụng thử nghiệm để kiểm chứng khả năng triển khai và đánh giá thực nghiệm để đối chiếu với các lập luận lý thuyết. Phần sau của chương tập trung vào giới hạn của nghiên cứu và những hướng cần tiếp tục làm sâu hơn nếu muốn đưa kết quả này tiến thêm một bước.

\section{Kết luận chung}

MLS \cite{rfc9420} cung cấp bảo mật mạnh mẽ cho quá trình giao tiếp nhóm, nhưng giao thức này mặc định dựa vào một Delivery Service tập trung để áp đặt thứ tự Commit. Khi đưa MLS vào mạng ngang hàng, không còn thực thể trung tâm nào đảm nhiệm vai trò đó, và hệ quả là các nút dễ rơi vào tình trạng phân nhánh mật mã không thể tự khôi phục. Đồ án này đã giải quyết bài toán trên bằng cách đề xuất một lớp điều phối phi tập trung đặt bao quanh MLS, đồng thời xây dựng một ứng dụng desktop thử nghiệm để kiểm chứng tính khả thi của hướng tiếp cận.

Về mặt thiết kế, đóng góp cốt lõi là việc phối hợp bốn cơ chế thành một giao thức điều phối nhất quán: Single-Writer theo epoch để chỉ một Token Holder được quyền tạo Commit, giảm xung đột ngay từ đầu; kiểm tra epoch để loại bỏ thông điệp cũ trước khi đưa xuống lớp mật mã; phát hiện phân nhánh và Fork Healing để đưa các nút hội tụ trạng thái sau phân vùng mạng; và đồng hồ logic lai HLC để ổn định thứ tự hiển thị thông điệp ứng dụng độc lập với thứ tự mật mã. Điểm quan trọng là toàn bộ thiết kế không sửa đổi lõi mật mã của MLS mà chỉ bổ sung lớp điều phối ở phía ngoài, qua đó giữ nguyên các bảo đảm bảo mật mà RFC 9420 cung cấp.

Một đóng góp đáng chú ý trong Fork Healing là việc dùng External Proposal kết hợp Token Holder thay vì External Commit. Khi nhiều nút cùng cần gia nhập lại sau phân vùng, External Proposal cho phép Token Holder gom tất cả đề nghị vào một Commit duy nhất, thay vì để mỗi nút tự tạo Commit cạnh tranh. So với các thuật toán đồng thuận truyền thống như Raft hay BFT yêu cầu trao đổi chéo $\mathcal{O}(N^2)$ thông điệp để chốt trạng thái, giao thức đề xuất chỉ cần Token Holder phát sóng một Commit, giữ độ phức tạp truyền thông ở mức $\mathcal{O}(1)$. Hơn nữa, các phân vùng nhỏ vẫn tiếp tục hoạt động độc lập khi mạng bị chia cắt thay vì đóng băng chờ quorum, phù hợp với đặc thù ứng dụng trò chuyện phi tập trung.

Về mặt triển khai, đồ án đã xây dựng một ứng dụng desktop tích hợp đầy đủ mạng P2P libp2p, cơ sở dữ liệu cục bộ SQLite và lõi mật mã MLS/OpenMLS qua kiến trúc sidecar stateless bằng Rust. Các luồng từ khởi tạo định danh, cấp bundle, quản trị tổ chức, tạo nhóm đến trao đổi tin nhắn được triển khai trọn vẹn, qua đó chứng minh rằng thiết kế giao thức không chỉ dừng ở mô tả lý thuyết mà có thể đi vào một hệ thống chạy được.

Về mặt đánh giá thực nghiệm, các kết quả đo đạc đã đối chiếu tích cực với các lập luận lý thuyết. Chaos test với cụm năm nút bị phân vùng lặp lại cho thấy các nút hội tụ về cùng epoch và cùng mã băm cây sau khi mạng được nối lại, không cần Delivery Service trung tâm. Thí nghiệm về cơ chế gom đề xuất cho thấy chiến lược Single-Writer giữ số Commit luôn bằng 1 và tỷ lệ thành công đạt 1,00, trong khi xử lý ngay tạo ra N Commit cạnh tranh với tỷ lệ thành công giảm dần khi số nút tăng. Benchmark mật mã cho thấy thời gian mã hóa MLS tăng theo $\mathcal{O}(\log N)$ trong khi mã hóa từng cặp tăng tuyến tính $\mathcal{O}(N)$ --- ở nhóm 4096 thành viên, MLS chỉ tốn 78,95 ms so với 955,18 ms của mã hóa từng cặp, chênh lệch hơn 12 lần. Tỷ trọng chi phí điều phối so với tổng chi phí cục bộ chỉ chiếm dưới 4\%, cho thấy lớp điều phối bổ sung không tạo gánh nặng đáng kể. Cuối cùng, kiểm thử với thành phần MLS thật xác nhận thuộc tính Forward Secrecy được giữ nguyên sau khi loại thành viên, chứng tỏ lớp điều phối không phá vỡ các bảo đảm mật mã của MLS.

\section{Giới hạn của đồ án}

Bên cạnh các kết quả đã đạt được, đồ án vẫn còn những giới hạn cần được nhìn nhận thẳng thắn. Trước hết, phần phân tích lý thuyết trong đồ án mới dừng ở mức lập luận theo bất biến và theo giả định vận hành chính, chưa đi tới kiểm chứng hình thức đầy đủ cho toàn bộ giao thức. Vì vậy, các kết luận của đồ án cần được hiểu trong phạm vi giả định đã nêu, không nên suy rộng thành một bảo đảm tuyệt đối cho mọi môi trường phân tán.

Tiếp theo, phần đánh giá thực nghiệm hiện nay vẫn chủ yếu dựa trên mạng giả lập và các bài kiểm thử có kiểm soát. Cách làm này phù hợp với phạm vi một đồ án tốt nghiệp vì nó cho phép cô lập các cơ chế cần kiểm tra. Tuy nhiên, nó chưa phản ánh hết độ phức tạp của mạng libp2p thực tế, đặc biệt khi có độ trễ lớn, thay đổi kết nối liên tục hoặc số lượng thiết bị tăng mạnh.

Một giới hạn khác đến từ chính ứng dụng thử nghiệm hiện tại. Trong nhiều thao tác, hệ thống vẫn trao đổi toàn bộ trạng thái nhóm giữa phần điều phối và phần MLS. Cách làm này giúp cấu trúc triển khai rõ ràng, nhưng tạo thêm chi phí khi quy mô nhóm tăng. Vì vậy, ứng dụng này hiện nay phù hợp với vai trò kiểm chứng khả thi, chưa phải là một phương án tối ưu cho triển khai quy mô lớn.

Thêm vào đó, cơ chế tự mã hóa và phát lại tin nhắn vẫn tồn tại điểm yếu liên quan đến bài toán trải nghiệm người dùng. Nếu một thành viên gửi tin nhắn trong lúc mạng phân mảnh nhưng lại ngoại tuyến trước khi mạng hồi phục và hợp nhất nhánh, hệ thống sẽ tạm thời mất đi thông điệp đó cho đến khi người gửi trực tuyến trở lại, do không ai khác có được khóa bí mật để mã hóa bản rõ thay cho họ.

Một giới hạn khác liên quan đến độ sâu phân kỳ. Thí nghiệm cho thấy thời gian healing tăng nhanh khi hai nhánh đã lệch nhiều epoch, từ 1374 ms ở độ sâu 10 lên 10370 ms ở độ sâu 100. Mặc dù External Proposal giữ số Commits luôn bằng 1, phần chi phí nằm ở phía nút nhánh thua khi phải xử lý Welcome và khôi phục trạng thái qua nhiều epoch. Đây là điểm cần tối ưu thêm nếu muốn hệ thống phục hồi nhanh hơn sau phân vùng kéo dài.

Cuối cùng, cần nhấn mạnh rằng ứng dụng desktop trong đồ án chỉ đóng vai trò ứng dụng thử nghiệm. Nó chứng minh rằng các cơ chế được đề xuất có thể đi vào luồng sử dụng thực tế, nhưng chưa đủ để xem như một sản phẩm hoàn thiện. Đóng góp chính của đồ án vẫn là lớp điều phối phi tập trung cho MLS trên mạng ngang hàng, còn ứng dụng chỉ là phương tiện triển khai và kiểm chứng hướng tiếp cận đó.

\section{Hướng phát triển}

Hướng phát triển đầu tiên là làm chặt hơn phần cơ sở lý thuyết. Cụ thể, các bất biến như tính an toàn của cơ chế Single-Writer, tính đơn điệu của epoch và điều kiện đạt hội tụ trạng thái mật mã sau Fork Healing (Định nghĩa~\ref{def:hoi_tu_trang_thai}) cần được mô hình hóa và kiểm chứng hình thức đầy đủ hơn. Nếu làm được điều này, giá trị học thuật của hướng nghiên cứu sẽ được nâng lên rõ rệt.

Hướng phát triển thứ hai là mở rộng thực nghiệm trên môi trường gần với triển khai thật hơn. Các thí nghiệm nên được chạy trên nhiều tiến trình, nhiều máy và trong các điều kiện mạng đa dạng hơn, chẳng hạn độ trễ cao, mất gói, thay đổi thành viên liên tục hoặc phân vùng mạng kéo dài. Khi đó, đồ án có thể cung cấp được bức tranh đầy đủ hơn về độ ổn định và khả năng chịu tải của hệ thống.

Hướng phát triển thứ ba là tối ưu hóa khâu truyền thông giữa phần điều phối và phần MLS. Việc giảm chi phí trao đổi trạng thái, tổ chức lại cách lưu trữ hoặc áp dụng các kỹ thuật snapshot hiệu quả hơn có thể giúp hệ thống giữ được kiến trúc tách lớp hiện tại nhưng hoạt động tốt hơn khi quy mô nhóm tăng. Đặc biệt, phần khôi phục trạng thái khi độ sâu phân kỳ lớn cần được tối ưu để giảm thời gian healing, chẳng hạn qua việc tăng tốc xử lý Welcome hoặc áp dụng cơ chế checkpoint định kỳ để giới hạn số epoch cần khôi phục.

Hướng phát triển cuối cùng là tiếp tục hoàn thiện ứng dụng thử nghiệm theo hướng phục vụ tốt hơn cho kiểm chứng và cho trình diễn. Các công cụ quan sát trạng thái, kiểm soát thay đổi thành viên và theo dõi quá trình hội tụ nếu được làm rõ hơn sẽ giúp việc đánh giá hệ thống thuận lợi hơn ở các nghiên cứu tiếp theo.

Kết chương. Chương này đã tổng hợp lại những đóng góp chính, giới hạn và hướng phát triển của đồ án. Có thể thấy kết quả quan trọng nhất của đồ án không nằm ở một ứng dụng đơn lẻ, mà nằm ở việc đề xuất được một lớp điều phối phi tập trung cho MLS trên mạng ngang hàng, rồi kiểm chứng được tính khả thi của hướng tiếp cận đó bằng cả phân tích lý thuyết, thực nghiệm và ứng dụng thử nghiệm chạy được.

Những phần tài liệu tham khảo và phụ lục ở sau chương này đóng vai trò đối chiếu nguồn và bổ sung minh họa cho các nội dung đã trình bày. Chúng giúp người đọc lần lại cơ sở học thuật của đồ án, đồng thời quan sát rõ hơn các màn hình và kịch bản của ứng dụng thử nghiệm đã được dùng trong phần kiểm chứng.

\end{document}
